\documentclass[12pt]{article}

% -----------------------
% Paquetes de idioma y formato
% -----------------------
\usepackage[utf8]{inputenc}
\usepackage[T1]{fontenc}
\usepackage[spanish]{babel}
\usepackage[margin=2.5cm]{geometry}
\usepackage{graphicx}
\usepackage{url}
\usepackage{microtype}   % Mejora tipográfica
\usepackage{booktabs}    % Tablas más profesionales
\usepackage{longtable}
\usepackage{array}
\usepackage{tabularx}
\usepackage{caption}     % Control de caption
\usepackage{float}       % [H] placement

\usepackage{amsmath}
\usepackage{amsfonts}
\usepackage{amssymb}
\usepackage{tikz}
\usepackage{algorithmicx}

\usepackage{natbib}
\usepackage{forest}
\usepackage{tikz}
\usetikzlibrary{trees}

% -----------------------
% Tipografía y espaciado
% -----------------------
\usepackage{mathptmx}
\usepackage{setspace}
\onehalfspacing
\setlength{\parskip}{8pt}

% -----------------------
% Encabezados
% -----------------------
\usepackage{fancyhdr}
\setlength{\headheight}{14.5pt}
\pagestyle{fancy}
\fancyhf{}
\renewcommand{\headrulewidth}{0pt}
\rhead{\thepage}

% -----------------------
% Código fuente (listings)
% -----------------------
\usepackage{listings}
\usepackage{xcolor}

% -----------------------
% Metadatos PDF (opcional)
% -----------------------
\usepackage[
    hidelinks,
    pdfauthor={Morales García, J. M., & Zamora Flores, D.},
    pdftitle={ADA_2.4_Determinación_de_umbral_MoralesG.ZamoraF}
]{hyperref}

% -----------------------
% Documento
% -----------------------
\begin{document}

% Portada
\begin{titlepage}
\begin{center}

\begin{minipage}{0.45\textwidth}
    \centering
    \includegraphics[width=0.8\textwidth]{Portada/Logo_IPN.png}
\end{minipage}
\hfill
\begin{minipage}{0.45\textwidth}
    \centering
    \includegraphics[width=0.6\textwidth]{Portada/Logo_UPIIT.png}
\end{minipage}

\vspace{1.5cm}

{\Large\bfseries Tarea 2.4\par}
\vspace{0.5cm}
{\LARGE\bfseries Determinación de Umbral\par}

\vspace{1.5cm}

{\large Instituto Politécnico Nacional\par}
{\large Unidad Profesional Interdisciplinaria de Ingeniería
Campus Tlaxcala\par}

\vspace{1.8cm}

\textbf{Integrante:}\\
Morales García, José Miguel \\
Zamora Flores, Daniel

\vspace{0.5cm}

\textbf{Curso:}\\
D301 -- Análisis y Diseño de Algoritmos

\vspace{0.5cm}

\textbf{Profesor:}\\
Escobar Juárez, Esaú Eliezer

\vfill

{\large 14 de octubre de 2026\par}

\end{center}
\end{titlepage}

\section{Introducción}
| Paso | Qué harás                                              | Programa que utilizarás                                            | Resultado que debes obtener                                   |
| ---- | ------------------------------------------------------ | ------------------------------------------------------------------ | ------------------------------------------------------------- |
| 1    | Generar puntos aleatorios.                             | Algoritmo de pares más cercanos.                                   | Un arreglo de puntos para las pruebas.                        |
| 2    | Medir fuerza bruta y determinar \(h(n)\).              | Algoritmo de pares más cercanos y programa de medición de tiempos. | Tabla de tiempos y análisis de su crecimiento.                |
| 3    | Medir la sección de combinación y determinar \(g(n)\). | Algoritmo de pares más cercanos y programa de medición de tiempos. | Tabla de tiempos de combinación y análisis de su crecimiento. |
| 4    | Medir los tiempos para \(n=10,20,\ldots,1000\).        | Algoritmo de pares más cercanos modificado.                        | Tabla comparativa de fuerza bruta y divide y vencerás.        |
| 5    | Determinar el umbral \(n_0\).                          | Resultados experimentales y análisis de los algoritmos.            | Un umbral justificado.                                        |
| 6    | Comparar los algoritmos y elaborar la gráfica.         | Algoritmo original modificado para incorporar el umbral.           | Gráfica y conclusiones sobre el umbral.                       |

\section{Generar puntos aleatorios}


% -----------------------
%Referencias
% -----------------------
\newpage
%\addcontentsline{toc}{section}{Referencias}
%\bibliographystyle{apalike}
%\bibliography{referencias}

\end{document}